\documentclass{situationist-preprint}

\title{Kant on Free Beauty}
\msid{AES-11-b}
\msfamily{AES-11}
\msdiscipline{Aesthetics}
\mstype{Companion manuscript}
\mskeywords{Kant; free beauty; dependent beauty; formalism; natural beauty; scientific cognitivism}

\begin{document}
\maketitle

\begin{abstract}
The formalist reply to scientific cognitivism about nature holds that much natural beauty is category-independent: the color of a sunset or the form of a flower can be appreciated without knowing what the object is. Kant's distinction between free and dependent beauty is the classical statement of this reply. We examine it and argue that it supports less than the formalist needs. Free beauty is real, and its appreciation requires no knowledge of kinds; but Kant himself classes the beauty of natural kinds such as animals and human beings as dependent, and the free beauties he lists are a narrow class of patterns and surfaces. The appreciation of natural beauty as a whole therefore still requires scientific rather than art-historical knowledge, as AES-11-a holds.
\end{abstract}

\section{Kant's distinction}

In §16 of the \emph{Critique of the Power of Judgment}, Kant distinguishes free beauty (\emph{pulchritudo vaga}), which presupposes no concept of what the object ought to be, from dependent or adherent beauty (\emph{pulchritudo adhaerens}), which presupposes such a concept and the perfection of the object in accordance with it \citep[§16]{Kant1790}. His examples of free beauty are flowers, many birds, a number of marine creatures, designs \emph{\`a la grecque} and foliage on wallpaper. Of the flower he remarks that hardly anyone but the botanist knows what it is meant to be, and that even the botanist, recognizing it as the reproductive organ of the plant, takes no account of this end when he judges it by taste. His examples of dependent beauty are the beauty of a human being, of a horse and of a building.

The formalist reads this as a precedent. If the botanist sets aside his knowledge in judging the flower beautiful, scientific knowledge is not required for the appreciation of natural beauty; and Objection~3 of AES-11-a makes exactly this claim.

\section{What the distinction supports}

\begin{proposition}\label{p:free}
Kant's account of free beauty establishes that some natural beauty is category-independent, and does not establish that natural beauty in general is.
\end{proposition}

\begin{proof}[Argument]
The first part is granted: a judgment of free beauty, on Kant's account, concerns form alone, and a pleasing pattern of color or symmetry can be appreciated by someone who does not know what the object is. The second part follows from Kant's own classification. The beauty of a horse is dependent: it presupposes a concept of what a horse is to be. The same holds, by Kant's reasoning, of animals generally, whenever they are judged as the animals they are. But to judge an animal as the kind of animal it is, in its correct category, is precisely what Principle~P3 and Lemma~3.2 of AES-11-a require, and that category is fixed by natural history. The free beauties Kant lists are those of surfaces and patterns considered apart from their kind. A formalism confined to them covers only part of natural beauty.
\end{proof}

Two further points narrow the precedent. First, the botanist sets aside the end of the flower, its function; he does not thereby perceive it outside every category. The perception of its form as graceful or delicate still depends on comparison with other flowers, and the categories of comparison are natural kinds \citep{Walton1970}. Second, Kant's free beauty is defined by the absence of a concept of purpose, not of a concept of kind. Scientific cognitivism does not require that natural objects be appreciated as purposive: it requires that they be perceived in the categories that explain their features, which include their origins and the processes that formed them \citep{Carlson2000,Budd2002}.

\begin{objection}
If free beauty is category-independent, scientific knowledge is not required for it, and the thesis of AES-11-a is false for this part of nature.
\end{objection}

\begin{reply}
The thesis concerns the appreciation of natural beauty as such, and AES-11-a concedes (in its reply to Objection~3) that some formal properties are category-independent. What it denies is that an aesthetic of nature can be built from them. Kant's own division agrees: the beauty that attaches to natural things as the things they are is dependent beauty.
\end{reply}

\section{Conclusion}

Kant gives the formalist a genuine phenomenon, and his division of beauty marks its limits. Free beauty in nature is the beauty of patterns; the beauty of natural things as the things they are depends on concepts of their kinds, and those concepts belong to natural history.

\begin{thebibliography}{9}\small
\bibitem[Budd(2002)]{Budd2002} Budd, M. (2002). \emph{The Aesthetic Appreciation of Nature: Essays on the Aesthetics of Nature}. Clarendon Press.
\bibitem[Carlson(2000)]{Carlson2000} Carlson, A. (2000). \emph{Aesthetics and the Environment: The Appreciation of Nature, Art and Architecture}. Routledge.
\bibitem[Kant(1790)]{Kant1790} Kant, I. (1790). \emph{Critique of the Power of Judgment}. Trans. P. Guyer and E. Matthews. Cambridge University Press, 2000.
\bibitem[Walton(1970)]{Walton1970} Walton, K.~L. (1970). Categories of art. \emph{The Philosophical Review}, 79(3), 334--367.
\end{thebibliography}

\end{document}
